\begin{pastexam}[label={px:gw3-108-ds-04}]{關務三等資訊處理 108 年 第四題}
（申論；20 分）

請將所給定數字藉由所指定雜湊函數依序置入雜湊表
\begin{center}
\begin{tabular}{|c|*{6}{p{7mm}<{\centering}|}}
\hline
A[.] & 0 & 1 & 2 & 3 & 4 & \dots \\
\hline
槽 1 & & & & & & \\
\hline
槽 2 & & & & & & \\
\hline
\end{tabular}
\end{center}
\begin{enumerate}[leftmargin=2.2em]
\item[㈠] 若雜湊函數為 H(k) = k mod 11，並以線性探測（linear probing）解決溢位（overflow）問題，請顯示將
15, 23, -12, 3, -8, 8, 9, 11, -3, -5, 14, 10, 25, 12, 0, 21 依序置入 11 桶（buckets）x 2 槽（slots）雜湊表的
最終結果。（10 分）
\item[㈡] 若雜湊函數為 H(k) = k mod 7，並以平方探測（quadratic probing）解決溢位（overflow）問題，請顯示將
15, 23, -12, 3, -8, 8, 9, 11, -3, -5, 14, 10, 25, 12 依序置入 7 桶（buckets）x 2 槽（slots）雜湊表的最終結果。
（10 分）
\end{enumerate}
\end{pastexam}

\begin{pastsoln}
\textbf{約定。} 每個桶有 2 個槽，桶裡還有空槽就放進去（先放槽 1 再放槽 2），兩槽都滿了才算溢位，改探測下一個
桶（資結書第~12.3~節）。 負數的 mod 取數學上的餘數，落在 0 到 $m-1$：例如 $-12 \bmod 11 = 10$（因為
$-12 = -2 \times 11 + 10$）。 C 語言的 \texttt{\%} 對負數會得到負的餘數（\texttt{-12 \% 11} 是 \texttt{-1}），
不能直接當桶號，要再加上 $m$。

\textbf{㈠ $H(k) = k \bmod 11$，線性探測。} 溢位時依序試 $H(k)+1, H(k)+2, \dots$（mod 11）。
\begin{center}
\small
\begin{tabular}{rclrclrcl}
\toprule
$k$ & $H(k)$ & 放入 & $k$ & $H(k)$ & 放入 & $k$ & $H(k)$ & 放入 \\
\midrule
15 & 4 & 4 & 9 & 9 & 9 & 10 & 10 & 10 \\
23 & 1 & 1 & 11 & 0 & 0 & 25 & 3 & 3 滿、4 滿 → 5 \\
$-12$ & 10 & 10 & $-3$ & 8 & 8 & 12 & 1 & 1 \\
3 & 3 & 3 & $-5$ & 6 & 6 & 0 & 0 & 0 \\
$-8$ & 3 & 3 & 14 & 3 & 3 滿 → 4 & 21 & 10 & 10 滿、0 滿、1 滿 → 2 \\
8 & 8 & 8 & & & & & & \\
\bottomrule
\end{tabular}
\end{center}
說明三個溢位的鍵：14 到桶 3 時已有 3、$-8$ 兩個，改放桶 4（和 15 同桶）；25 時桶 3、桶 4 都滿，放桶 5；
21 時桶 10 已有 $-12$、10，往後繞回桶 0（11、0）、桶 1（23、12）都滿，放桶 2。 最終結果：
\begin{center}
\begin{tabular}{|c|*{11}{c|}}
\hline
A[.] & 0 & 1 & 2 & 3 & 4 & 5 & 6 & 7 & 8 & 9 & 10 \\
\hline
槽 1 & 11 & 23 & 21 & 3 & 15 & 25 & $-5$ & & 8 & 9 & $-12$ \\
\hline
槽 2 & 0 & 12 & & $-8$ & 14 & & & & $-3$ & & 10 \\
\hline
\end{tabular}
\end{center}

\textbf{㈡ $H(k) = k \bmod 7$，平方探測。} 第 $i$ 次探測桶 $(H(k) + i^2) \bmod 7$，$i = 0, 1, 2, \dots$，也就是
$H, H+1, H+4, H+9, \dots$。
\begin{center}
\small
\begin{tabular}{rcl}
\toprule
$k$ & $H(k)$ & 探測與結果 \\
\midrule
15 & 1 & 桶 1 \\
23 & 2 & 桶 2 \\
$-12$ & 2 & 桶 2（滿了） \\
3 & 3 & 桶 3 \\
$-8$ & 6 & 桶 6 \\
8 & 1 & 桶 1（滿了） \\
9 & 2 & 桶 2 滿；$i = 1$：桶 3 \\
11 & 4 & 桶 4 \\
$-3$ & 4 & 桶 4（滿了） \\
$-5$ & 2 & 桶 2 滿；$i=1$：桶 3 滿；$i=2$：$(2+4) \bmod 7 = 6$，桶 6 \\
14 & 0 & 桶 0 \\
10 & 3 & 桶 3 滿；$i=1$：桶 4 滿；$i=2$：$(3+4) \bmod 7 = 0$，桶 0 \\
25 & 4 & 桶 4 滿；$i=1$：桶 5 \\
12 & 5 & 桶 5（滿了） \\
\bottomrule
\end{tabular}
\end{center}
14 個數剛好填滿 7 桶 $\times$ 2 槽。 最終結果：
\begin{center}
\begin{tabular}{|c|*{7}{c|}}
\hline
A[.] & 0 & 1 & 2 & 3 & 4 & 5 & 6 \\
\hline
槽 1 & 14 & 15 & 23 & 3 & 11 & 25 & $-8$ \\
\hline
槽 2 & 10 & 8 & $-12$ & 9 & $-3$ & 12 & $-5$ \\
\hline
\end{tabular}
\end{center}
平方探測不保證走遍所有桶：$i^2 \bmod 7$ 只會是 0、1、2、4，從桶 $H$ 出發只到得了 $H$、$H+1$、$H+2$、$H+4$
四個桶（資結書定理~12.3）。 本題每個鍵都在這四個桶裡找到空槽，所以全部放得進去。 若用 $H \pm i^2$ 交替
探測的版本，這組資料的最終結果也一樣。
\end{pastsoln}
